\section*{Chapter \ref{ch:qm:girsanov-and-changes-of-numeraire} --- Girsanov and Changes of Numeraire}

\begin{solution}{exo:qm:girsanov-and-changes-of-numeraire:1}
$W_T = W^{\mathbb Q}_T + \gamma T$ with $W^{\mathbb Q}$ a $\mathbb Q$-Brownian motion: $W_T \sim \mathcal N(2, 4)$ under $\mathbb Q$.
\end{solution}

\begin{solution}{exo:qm:girsanov-and-changes-of-numeraire:2}
Drift 4\%; $\gamma = -(\mu - r)/\sigma = -0.05/0.25 = -0.2$, the negative of the Sharpe ratio.
\end{solution}

\begin{solution}{exo:qm:girsanov-and-changes-of-numeraire:3}
$P(0,5)\,\E^5[X] = 0.8343 \times 2.4 = 2.002$.
\end{solution}

\begin{solution}{exo:qm:girsanov-and-changes-of-numeraire:4}
Under $\mathbb Q^T$ the drift of $r$ loses $\sigma^2B(T - t) > 0$: states with high rates, in which the bond
is cheap, get less weight. $B(5) = 1.8358$ and $\sigma^2B^2/2 = 0.0001 \times 3.370/2 = 1.7$ bp, the gap between
3.918\% and 3.901\%.
\end{solution}

\begin{solution}{exo:qm:girsanov-and-changes-of-numeraire:5}
$S_t = (P(t,T_0) - P(t,T_n))/A_t$ is a traded portfolio divided by the \omterm{def:qm:girsanov-and-changes-of-numeraire:numeraire}{numeraire} $A_t$, hence a
$\mathbb Q^A$-martingale. A receiver swaption pays $A_{T_0}(K - S_{T_0})^+$ at $T_0$, so its value is $A_0\,\E^A[(K -
S_{T_0})^+]$.
\end{solution}

\begin{solution}{exo:qm:girsanov-and-changes-of-numeraire:6}
Under $\mathbb Q^{S^2}$ the ratio $S^1/S^2$ is a \omterm{def:qm:probability-at-speed:martingale}{martingale} with volatility $\sqrt{0.2^2 + 0.2^2 - 2 \times 0.5
\times 0.04} = 0.2$, and the value is $S^2_0\,\E^{S^2}[(S^1_T/S^2_T - 1)^+] = 100(\Phi(0.1) - \Phi(-0.1)) = 7.97$:
Margrabe's formula, with no interest rate in it.
\end{solution}

\begin{solution}{exo:qm:girsanov-and-changes-of-numeraire:7}
Let $X = \mu t + \sigma W$. Under $d\mathbb Q/d\P = e^{-(\mu/\sigma)W_T - \mu^2T/2\sigma^2}$, $X/\sigma$ is a standard \omterm{def:qm:brownian-motion:bm}{Brownian
motion}, and $d\P/d\mathbb Q = e^{(\mu/\sigma^2)X_T - \mu^2T/2\sigma^2}$ depends on the path only through $X_T$. With $m =
\min_{t\le T}X_t$ and $b < 0$, the reflection principle gives under $\mathbb Q$ the joint law $\mathbb Q(m \le b, X_T \in dx)
= \mathbb Q(X_T \in 2b - dx)$ for $x \ge b$. Hence
\[
\P(m \le b) = \P(X_T \le b) + \int_b^\infty e^{\mu x/\sigma^2 - \mu^2T/2\sigma^2}\,\varphi_{\sigma\sqrt T}(2b - x)\,dx,
\]
where $\varphi_s$ is the $\mathcal N(0, s^2)$ density; completing the square, the integral is $e^{2\mu b/\sigma^2}\Phi((b +
\mu T)/\sigma\sqrt T)$, and $\P(X_T \le b) = \Phi((b - \mu T)/\sigma\sqrt T)$. For the stop of chapter~2 ($b = \ln
0.98$, $\mu = -0.045$, $\sigma = 0.30$, $T = 21/252$): 82.4\%.
\end{solution}

\begin{solution}{exo:qm:girsanov-and-changes-of-numeraire:8}
A price is an expectation under a \omterm{def:qm:probability-at-speed:martingale}{martingale} measure for the chosen \omterm{def:qm:girsanov-and-changes-of-numeraire:numeraire}{numeraire}, not under the
real-world measure: with the \omterm{def:qm:girsanov-and-changes-of-numeraire:numeraire}{money-market account} as \omterm{def:qm:girsanov-and-changes-of-numeraire:numeraire}{numeraire} the stock must drift at the short
rate. Using the forecast drift prices a different, arbitrageable claim; the forecast belongs in the
trading decision (buy if the market price is below the value one's view implies), not in the
pricing measure.
\end{solution}

\begin{solution}{pb:qm:girsanov-and-changes-of-numeraire:1}
\textbf{1.} $P(0,5) = 0.8343$, $P(0,5.25) = 0.8262$.
\textbf{2.} $(0.8343/0.8262 - 1)/0.25 = 3.92\%$.
\textbf{3.} 3.918\% and 3.901\%: under the \omterm{def:qm:girsanov-and-changes-of-numeraire:forward}{forward measure} the drift of $r$ is lower by $\sigma^2B(5 - t)$.
\textbf{4.} A put struck at $x = 1/(1 + \delta K) = 0.98888$ on $1 + \delta K = 1.01125$ zero-coupon bonds maturing
at 5.25, expiring at 5.
\textbf{5.} $\sigma_P = 0.01\sqrt{(1 - e^{-5})/1}\times B(0.25) = 0.234\%$.
\textbf{6.} 3.26 bp of notional: USD~32\,600.
\textbf{7.} 260 (weekly steps over five years) against one.
\textbf{8.} $3.34 \pm 0.05$ bp.
\textbf{9.} $3.32 \pm 0.05$ bp.
\textbf{10.} Yes: 1.6 and 1.2 standard errors from 3.26.
\textbf{11.} $(P(5,5)/P(0,5))/(B_5/B_0) = 1/(P(0,5)B_5)$, with $B_5 = e^{\int_0^5r}$ along the path.
\textbf{12.} $3.906\% \pm 0.003\%$, against $f(0,5) = 3.901\%$; the weights average 1.000.
\textbf{13.} 0.996: the discount factor's randomness is small next to the payoff's, so removing it
barely reduces the variance per path.
\textbf{14.} The cost per path: one normal instead of 260, so about 260 times less work for the same
error, and then the closed form with none.
\textbf{15.} No: $3.32 \pm 0.05$ bp with 208 steps a year, within noise of 3.34.
\textbf{16.} Each caplet under the \omterm{def:qm:girsanov-and-changes-of-numeraire:forward}{forward measure} of its own date, then sum; a joint simulation of all
dates needs one measure for all, the spot or terminal measure of One Quant Book~6.
\textbf{17.} The \omterm{def:qm:girsanov-and-changes-of-numeraire:annuity}{annuity measure}; the par swap rate is its \omterm{def:qm:probability-at-speed:martingale}{martingale}.
\textbf{18.} That every simulated price divided by the \omterm{def:qm:girsanov-and-changes-of-numeraire:numeraire}{numeraire} has a constant mean over dates (the
deflated \omterm{def:qm:probability-at-speed:martingale}{martingale} test), and that change-of-numeraire weights average one.
\textbf{19.} Named result: \emph{two measures, one price}: the caplet is worth 3.26 bp of notional
(USD~32\,600 on USD~100 million); risk-neutral and forward-measure simulations agree with it and have
the same standard error per path, but the \omterm{def:qm:girsanov-and-changes-of-numeraire:forward}{forward measure} needs one normal a path instead of 260,
and yields the closed form.
\textbf{20.} It changes the drift and the unit in which values are counted, never the prices and never
the volatility.
\end{solution}

\begin{solution}{iq:qm:girsanov-and-changes-of-numeraire:1}
It adds a drift: $W - \int\gamma\,dt$ is a \omterm{def:qm:brownian-motion:bm}{Brownian motion} under the new measure. It cannot change the \omterm{def:qm:brownian-motion:qv}{quadratic
variation}, which is visible on a single path, so volatilities are the same under all \omterm{def:qm:probability-at-speed:change}{equivalent
measures}.
\iqlookfor{drift yes, volatility no, and why.}
\end{solution}

\begin{solution}{iq:qm:girsanov-and-changes-of-numeraire:2}
The \omterm{def:qm:girsanov-and-changes-of-numeraire:emm}{risk-neutral measure} is a pricing device, not a forecast: it is the measure under which prices
counted in units of the \omterm{def:qm:girsanov-and-changes-of-numeraire:numeraire}{money-market account} are fair games. A stock held against a short
money-market position must then earn nothing on average in those units, so its drift is $r$; the real
expected return lives under $\P$.
\iqlookfor{measure as a unit of account, not a belief.}
\end{solution}

\begin{solution}{iq:qm:girsanov-and-changes-of-numeraire:3}
The \omterm{def:qm:girsanov-and-changes-of-numeraire:emm}{equivalent martingale measure} for the zero-coupon bond maturing at $T$. The simple forward rate for
$[S, T]$ is $(P(t,S)/P(t,T) - 1)/\delta$: a traded price divided by the \omterm{def:qm:girsanov-and-changes-of-numeraire:numeraire}{numeraire}, minus a constant, hence a
\omterm{def:qm:probability-at-speed:martingale}{martingale}.
\iqlookfor{\omterm{def:qm:girsanov-and-changes-of-numeraire:numeraire}{numeraire} $P(t,T)$ and the ratio argument.}
\end{solution}

\begin{solution}{iq:qm:girsanov-and-changes-of-numeraire:4}
With asset 2 as \omterm{def:qm:girsanov-and-changes-of-numeraire:numeraire}{numeraire} the payoff $(S^1_T - S^2_T)^+ = S^2_T(S^1_T/S^2_T - 1)^+$ becomes a call on the
\omterm{def:qm:probability-at-speed:martingale}{martingale} $S^1/S^2$ struck at one, with volatility $\sqrt{\sigma_1^2 + \sigma_2^2 - 2\rho\sigma_1\sigma_2}$: Margrabe's
formula, independent of rates.
\iqlookfor{choosing the \omterm{def:qm:girsanov-and-changes-of-numeraire:numeraire}{numeraire} that removes one asset.}
\end{solution}

\begin{solution}{iq:qm:girsanov-and-changes-of-numeraire:5}
Simulate traded instruments whose prices are known (zero-coupon bonds, forwards) and check that
their prices divided by the engine's \omterm{def:qm:girsanov-and-changes-of-numeraire:numeraire}{numeraire} have constant means over all dates within a few
standard errors; check put--call parity on the simulated payoffs; and price the same claim under two
\omterm{def:qm:girsanov-and-changes-of-numeraire:numeraire}{numeraires} with the reweighting.
\iqlookfor{\omterm{def:qm:probability-at-speed:martingale}{martingale} tests on instruments with known prices.}
\end{solution}

\begin{solution}{iq:qm:girsanov-and-changes-of-numeraire:6}
Plain sampling sees an exceedance once in 3.5 million draws. Sample $Y \sim \mathcal N(5, 1)$ instead and average
$e^{-5Y + 12.5}\mathbf 1_{\{Y > 5\}}$: 10\,000 draws give $2.79 \times 10^{-7}$ with a standard error of
$0.07 \times 10^{-7}$, against the exact $2.87 \times 10^{-7}$.
\iqlookfor{importance sampling by a mean shift, with the likelihood-ratio weight.}
\end{solution}
